\section{Declaración del Uso de IA}

En cumplimiento con las buenas prácticas académicas y de transparencia, se declara que para la elaboración de este trabajo se ha contado con el apoyo de herramientas de Inteligencia Artificial Generativa, concretamente Claude y Claude Code, de Anthropic. El uso de dicha tecnología se ha limitado a las siguientes tareas:

\begin{itemize}[leftmargin=1.5em]
  \item \textbf{Scaffolding inicial del proyecto.} Se ha utilizado Claude Code para generar la estructura base del proyecto, incluyendo la configuración del servidor Express, las rutas de la API REST y los componentes React iniciales. Este código fue revisado, ejecutado y modificado por los autores durante el desarrollo.
  \item \textbf{Diseño del modelo de datos.} Se ha empleado Claude como herramienta de consulta para contrastar decisiones de diseño sobre el esquema de documentos en CouchDB, la definición de las vistas MapReduce y la elección entre índices Mango y vistas para cada tipo de consulta.
  \item \textbf{Soporte técnico durante el desarrollo.} Se ha recurrido a la IA como apoyo en la resolución de problemas concretos surgidos durante la implementación, principalmente relacionados con la configuración de autenticación en CouchDB y la integración con la API de TMDB.
  \item \textbf{Revisión de texto.} Se ha realizado una revisión de la redacción de esta memoria con apoyo de IA para mejorar la fluidez y coherencia del texto. Este proceso se ha aplicado siempre sobre contenido escrito originalmente por los autores, manteniendo íntegras las decisiones técnicas, argumentaciones y conclusiones del trabajo.
\end{itemize}

El Cuadro~\ref{tab:uso-ia} resume las partes del trabajo donde se utilizó IA y con qué propósito.

\begin{table}[H]
  \centering
  \caption{Detalle del uso de IA por parte del trabajo.}
  \label{tab:uso-ia}
  \begin{tabularx}{\textwidth}{llX}
    \toprule
    \textbf{Parte del trabajo} & \textbf{Herramienta} & \textbf{Propósito} \\
    \midrule
    Elección del tema y stack   & Claude (chat)  & Evaluar opciones y decidir tecnologías. \\
    Generación del código base  & Claude Code    & Scaffolding inicial de los 30 ficheros del proyecto. \\
    Diseño del modelo CouchDB   & Claude (chat)  & Definir esquema de documentos y vistas MapReduce. \\
    Depuración de configuración & Claude (chat)  & Resolver problemas de autenticación con CouchDB. \\
    Redacción de la memoria     & Claude (chat)  & Apoyo en la estructura y revisión del texto. \\
    \bottomrule
  \end{tabularx}
\end{table}

\subsection{Prompts utilizados}

A continuación se recogen los prompts empleados en las conversaciones con Claude, formulados para obtener orientación y apoyo sin delegar la toma de decisiones ni la implementación final en la herramienta.

\subsubsection*{Elección del tema y del stack}
\begin{quote}
\textit{Estoy valorando usar CouchDB para un proyecto académico de catálogo personal de películas. Antes de decidirme, quiero entender qué ventajas e inconvenientes tiene frente a un enfoque relacional clásico en este contexto concreto, y qué stack web tendría sentido usar: qué tecnologías encajan bien con CouchDB, por qué, y qué sacrificaría con cada opción. No quiero que tomes la decisión por mí, sino que me des los argumentos para tomarla yo.}
\end{quote}

\subsubsection*{Scaffolding del proyecto}
\begin{quote}
\textit{Necesito organizar la estructura de carpetas y ficheros de una aplicación con frontend y backend separados antes de empezar a implementar. Dime qué estructura tiene sentido, qué responsabilidad debería tener cada parte y cómo separarlas bien desde el principio. No hace falta que desarrolles el código, solo que me orientes sobre cómo plantearlo.}
\end{quote}

\subsubsection*{Modelo de datos en CouchDB}
\begin{quote}
\textit{Estoy diseñando el modelo de datos de un catálogo de películas en CouchDB. Necesito decidir qué campos debería tener un documento de película y por qué tiene sentido representarlo así en un modelo documental. Ayúdame a razonar la estructura, señalando qué encaja bien en un documento JSON y qué habría que repensar, pero sin escribirme la solución entera.}
\end{quote}

\subsubsection*{Vistas MapReduce}
\begin{quote}
\textit{Quiero implementar vistas MapReduce en CouchDB para agregar datos del catálogo: cuántas películas hay por género, cuántas por estado y cuál es la valoración media por categoría. Explícame cómo enfocar este tipo de consultas en CouchDB, qué hace la función map y qué papel juega el reduce en cada caso, pero déjame a mí escribir la implementación final.}
\end{quote}

\subsubsection*{Depuración de errores}
\begin{quote}
\textit{Tengo un problema con CouchDB en mi proyecto: el servidor arranca pero no consigue conectar con la base de datos y devuelve un error de autenticación. Necesito entender qué puede estar fallando. Dime qué comprobaciones haría para diagnosticar esto: dónde mirar, qué revisar y en qué orden, sin limitarte a darme una solución cerrada.}
\end{quote}

\subsubsection*{Redacción de la memoria}
\begin{quote}
\textit{He redactado varios apartados de la memoria y quiero que suenen más claros y técnicamente correctos. Te paso el texto y necesito que me ayudes a reorganizarlo y reformularlo para que tenga un tono más académico y sea más fácil de leer. El contenido y las conclusiones son míos y no deben cambiar; solo quiero mejorar cómo están expresados.}
\end{quote}

\subsection{Autoría y revisión crítica}

Los autores declaran conocer y poder explicar en detalle todo el código y el contenido entregado. Las respuestas obtenidas de Claude se trataron en todo momento como material orientativo, revisado críticamente antes de incorporarse al proyecto. La selección del dominio, el diseño del modelo de datos, la implementación, la ejecución de pruebas, el análisis de resultados y la redacción definitiva de la memoria fueron responsabilidad de los autores.
